\section{Actual ordinary products have nonrecovering sources after their last mode through half}
This is a weaker-curvature theorem, not an ordinary half-log-concavity
theorem. The proof is an induction on the total number of non-one arm
parts. It uses established single-arm source shape, ordinary-product
unimodality, interval-product derivative resources, first-missing mass,
and the completed both-forward-nonincreasing first-half theorem. It uses
neither property C nor OLD A1. All comparisons are of actual coefficients.

For a polynomial $p$, zero-extend its coefficients and write
\[
 E_p(n)=2p_n-p_{n-1},\qquad
 T_p(n)=E_p(n+1)-E_p(n)=2p_{n+1}-3p_n+p_{n-1}.
\]
\begin{theorem}[Ordinary half-tail source nonrecovery]\label{thm:actual-half-tail}
Let $X,Y$ be normalized actual increasing nonempty arms, and put
$p=q_Xq_Y$, $D=\operatorname{wt}X+\operatorname{wt}Y$.
If $\nu$ is the LAST ordinary coefficient mode of $p$, then
\[
                        T_p(n)\le0\qquad(\nu<n,\ 2n\le D). \tag{W}
\]
No nonnegative-source hypothesis is required. In fact $E_p(n)\ge0$
through this whole half range independently of (W).
\end{theorem}

\paragraph{Actual inputs used in the proof.}
A normalized arm $Z=1^\alpha U$ has non-one count $h$, weight $w$,
ordinary polynomial $q_Z$, main interval product
$K_Z=[\alpha+1]\prod_{u\in U}[u]$ of degree $w-h$, and corrected source
\[
 e_Z=(2-z)q_Z-z^w+z^{w+1}.
\]
The established actual inputs are:
\begin{itemize}
\item $q_Z$ is positive on $[0,\operatorname{wt}Z]$, unimodal,
and every ordinary height interval has endpoint sum at most its weight;
\item $e_Z$ is nonnegative unimodal, positive exactly on
$[0,\operatorname{wt}Z+1-\kappa_Z]$, starts with coefficient two,
and every height interval has endpoint sum at most $\operatorname{wt}Z-1$;
\item $q_Z\succeq K_Z$ and
$(q_Z(j)-q_Z(j-1))^+\le K_Z(j)$ at every rank;
\item $2h\le\operatorname{wt}Z-1$ and
$\operatorname{wt}Z\ge\kappa_Z(\kappa_Z-1)/2$;
\item for a nonrectangular arm, its second ordinary height interval is
exactly $[1,\operatorname{wt}Z-2]$.
\end{itemize}
The main-box derivative bound is
\texttt{main\_box\_derivative\_envelope.tex}; the mass bound is
\texttt{first\_defect\_mass.tex}. The remaining shape inputs are the
established ordinary and corrected-source layer theorems.

\paragraph{A window identity and its crossing case.}
Suppose $p$ is unimodal, with first mode $m$ and last mode $\ell$,
and $w=[t]p$. On the ordinary tail of $w$, put $k=n+1-t$.
Then
\[
 T_w(n)=E_p(n+1)-E_p(k),\qquad
 \Delta w_{n+1}-\Delta w_n=\Delta p_{n+1}-\Delta p_k.   \tag{1}
\]
The first window mode is at least $m$. If $k\le\ell$, the leaving
input increment is nonnegative, and the entering input increment is
nonpositive because $n$ is after the window mode and therefore at least
the first old mode. Hence $\Delta w_{n+1}\le\Delta w_n$; since the window
is on its ordinary tail, $\Delta w_{n+1}\le0$ as well. Therefore
\[
 T_w(n)=\Delta w_{n+1}+(\Delta w_{n+1}-\Delta w_n)\le0. \tag{2}
\]
This includes all maximum plateaus and negative leaving indices.
Only the case $k>\ell$ needs a source comparison.

\begin{lemma}[Rectangle-factor base]
Statement (W) holds whenever at least one arm is rectangular.
\end{lemma}
\begin{proof}
Write $p=[A+1]q$ with $\deg q=b$.
If $A\ge b$, its maximum plateau is exactly $[b,A]$, so there is no
strict post-LAST-mode position through $(A+b)/2$.
If $A<b$, write $e=(2-z)q-z^b+z^{b+1}$. Its proved unimodality and
the causal averaging relation $2q_j=e_j+q_{j-1}$ below rank $b$ imply
that the source mode is no later than the first $q$ mode. Consequently
the raw source of $q$ is nonincreasing after its ordinary mode up to
rank $b-1$. The window argument (1)--(2) pays every required comparison
with $n+1<b$: in its noncrossing case both source endpoints lie in
this monotone prefix.

The only remaining possibility is $n+1=b$, which the half-range
condition permits only for $b-A\in\{1,2\}$. The exact identity
\[
                  (2-z)p=[A+1]e+z^b-z^{A+b+1}
\]
gives
$T_p(b-1)=1-e_{b-A-1}\le0$, since $e_0=2$ and $e_1\ge1$.
Two rectangular factors can also be handled directly by their
log-concave trapezoidal product.
\end{proof}

\paragraph{Induction setup and the NEW terminal-source edges.}
Assume both new arms are nonrectangular. Choose the NEW heavier arm,
of weight $a\ge b$, delete its last non-one part $t$, and call its
old weight $c=a-t$. Let $p$ be the old ordinary product, and $P$ the
new one. Then
\[
                          P=[t]p+z^a q_Y.              \tag{3}
\]
The total non-one count decreases by one, so (W) holds for $p$ by
induction. The first mode of $P$ is strictly below $a$.
At that mode $P$ agrees with $w=[t]p$, and $w\le P$ everywhere,
so $\max w=\max P$ and every maximum of $w$ is also a maximum of $P$.
In particular
$\operatorname{lastmode}w\le\operatorname{lastmode}P$.
Thus a strict ordinary-tail position of $P$ is on the window tail.

The source in (3) can affect a half-range coefficient triple only when
$n\ge a-1$, hence only for NEW weight gaps $a-b=0,1,2$.
Those edges are all covered by the proved theorem ``the entire
both-nonincreasing first half'' in
\texttt{agent\_bothdeclining\_first\_birth.tex}:
\[
\begin{array}{c|c|c}
a-b&n&j=a+b-n\\\hline
2&a-1&a-1=b+1\\
1&a-1&a\\
0&a-1,\ a&a+1,\ a
\end{array}
\]
At each listed forward edge both actual factors are nonincreasing,
and $2j\le a+b+2$. That theorem gives $T_P(n)\le0$ without a mode,
source-sign or OLD assumption. Henceforth
\[
                              n+1<a.                   \tag{4}
\]
Then $T_P(n)=T_w(n)$ and $k=n+1-t<c$.
By (2) a possible failure must have $k>\ell$, the LAST old mode.

\paragraph{An OLD rectangular arm created by deletion.}
The unchanged old opposite arm is nonrectangular. If the deleted arm
has become the rectangle $[c+1]$, the old product is $[c+1]q_Y$ and
its LAST mode is at least $c$: it strictly increases through rank $c$
when $c<b$, and has maximum plateau $[b,c]$ when $c\ge b$.
But (4) gives $k<c\le\ell$, so (2) pays the edge.
The remaining induction case therefore has BOTH old arms nonrectangular.

\paragraph{The OLD positive source prefix and the entering correction.}
Orient the OLD weights as $M\ge s$, write $d=M+s$, and let
$\kappa\ge3$ be the larger old arm's first missing weight. Put
\[
 R=M+1-\kappa,\quad
 C=e_Lq_S,\quad E=(2-z)p.
\]
Then
\[
                         E=C+z^M(1-z)q_S.              \tag{5}
\]
Thus $E_j=C_j\ge0$ below $M$.
Every interval-product component of $C$ has endpoint sum at most
$M+s-1=d-1$, so $C$ is nonincreasing from $\lfloor d/2\rfloor$ onwards.
Together with the old induction hypothesis, this proves
\[
 E_u\ge E_v\quad(\ell<u\le v<M).                       \tag{6}
\]
No global sign-prefix or OLD A1 assertion is used.

If $n\le M-2$, both source endpoints in (1) lie in this old prefix,
so (6) pays the edge. Suppose now $n\ge M-1$.

\begin{lemma}[The initial entering-source boundary]
The position $n=M-1$ is always paid.
\end{lemma}
\begin{proof}
Here the entering value is $E_M=C_M+1$.
If $\kappa\le s+2$ and $M>s$, pair the lowest source layer $[0,R]$
with the lowest ordinary layer $[0,s]$. Its trapezoid loses exactly
one at $M-1\to M$, because the leaving point $\kappa-2$ lies in
$[0,s]$. All other component differences are nonnegative, so
$E_M\le E_{M-1}\le E_k$.
If $M=s$, use the second ordinary layer $[1,s-2]$ instead: the leaving
point $\kappa-2$ belongs to it because $3\le\kappa\le M$.

If $\kappa>s+2$, then $C$ has support endpoint
$R+s=M+s+1-\kappa\le M-2$, so $E_M=1$.
The half-range condition gives $k=M-t\le s+2$.
Since $s\ge3$, this case has $\kappa\ge6$, and first-missing mass
gives $R\ge2$. The convolution $C$ is positive on all of $[0,R+s]$,
which contains $[0,s+2]$. Since $k>\ell\ge0$ and $k<M$,
$E_k=C_k\ge1=E_M$.
\end{proof}

\paragraph{All later entering-source positions.}
Write $n=M+r$, $r\ge0$. Equations (1) and (5) give
\[
 T_P(n)=C_{M+r+1}+q_S(r+1)-q_S(r)-E_k.                 \tag{7}
\]
The new half-range condition gives
\[
                      k\le k_*:=s-r+1.                 \tag{8}
\]
Every comparison of a component at $M-1$ with $M+r+1$ has nonnegative
drop, since the two indices sum to $2M+r\ge d-1$.

\begin{lemma}[Late entering-source payment]
If $r+3\ge\kappa$, then (7) is nonpositive.
\end{lemma}
\begin{proof}
For each $q_S$ height interval $[A,B]$ starting at $A=r+1$, pair it
with the lowest source layer $[0,R]$. At $M-1$ this convolution counts
the entire interval, because $B\le s-r-1\le M-1$ and
$M-1-R=\kappa-2\le A$. At $M+r+1$ it misses at least $A$, since
$M+r+1-R=r+\kappa\ge A+1$.
Thus the total drop is at least the number of such starting layers,
which is at least $q_S(r+1)-q_S(r)$. Therefore
$E_{M+r+1}\le E_{M-1}\le E_k$, by (6).
\end{proof}

Only $0\le r\le\kappa-4$ remains. Let $h_S$ be the old short arm's
non-one count and $m_S=s-h_S$ the degree of its main interval product.

\begin{lemma}[A short main-box payment]
If $m_S\ge\kappa+r-1$, then (7) is nonpositive.
\end{lemma}
\begin{proof}
The lowest source layer contributes exactly
\[
 \sum_{u=\kappa-2}^{\kappa+r-1}q_S(u)
                  -\sum_{u=M}^{M+r+1}q_S(u)             \tag{9}
\]
to $C_{M-1}-C_{M+r+1}$. The second sum is zero if $M>s$, and one if
$M=s$. Let $K_S$ be the short arm's main interval product. Its symmetry
and unimodality, together with $r+1\le\kappa-2$ and
$\kappa+r-1\le m_S$, give
\[
 q_S(\kappa-2)\ge K_S(\kappa-2)\ge K_S(r+1)
                         \ge(q_S(r+1)-q_S(r))^+.
\]
There is at least one further positive coefficient in (9), namely
$q_S(\kappa-1)\ge1$, since $\kappa-1\le m_S\le s$.
It pays the possible terminal penalty one. Other component drops are
nonnegative, so again $E_{M+r+1}\le E_{M-1}\le E_k$.
\end{proof}

\begin{lemma}[A virtual central head payment]
Suppose $k_*<M$ and $\kappa\le M-s+2r+1$. Then (7) is nonpositive.
The number $M-s+2r$ is only a comparison width, not a claimed legal arm part.
\end{lemma}
\begin{proof}
Compare $C_{k_*}$ with $C_{M+r+1}$. Their indices sum to $d+2$,
so every component has nonnegative drop. For a $q_S$ height interval
beginning at $A=r+1$, its right endpoint obeys
$B\le s-r-1=k_*-2$. At the head $k_*$ the lowest source layer counts
all of $[A,B]$, since
\[
                   k_*-R=s-r+1-(M+1-\kappa)\le A.
\]
At the tail it misses $A$ as in the late payment.
Thus $E_{M+r+1}\le C_{k_*}=E_{k_*}\le E_k$ by (6) and (8).
\end{proof}

\paragraph{The two resources cover every nonexceptional head.}
Suppose $k_*<M$. If both preceding payments failed, then
\[
 M\le s+\kappa-2r-2,\qquad m_S\le\kappa+r-2.
\]
Normalization $2m_S\ge s+1$ gives $s\le2\kappa+2r-5$, and hence
$M\le3\kappa-7$. This contradicts first-missing mass:
\[
 \kappa(\kappa-1)/2>3\kappa-7
       \quad(\kappa\ge3),
\]
because the doubled gap is
$(\kappa-3)(\kappa-4)+2>0$ for every integer $\kappa\ge3$.

\paragraph{The three exceptional heads are also covered.}
The condition $k_*=s-r+1\ge M\ge s$ has only three possibilities.
\begin{itemize}
\item $M=s,r=0$: mass gives $s\ge2\kappa-3$, so
$2m_S\ge s+1$ implies $m_S\ge\kappa-1$ and the main-box payment applies.
\item $M=s+1,r=0$: mass gives $M\ge2\kappa-3$, and
$2m_S\ge s+1=M$ implies the same integer bound $m_S\ge\kappa-1$.
\item $M=s,r=1$: for $\kappa\le4$ the late payment already applies;
for $\kappa\ge5$, mass gives $s\ge2\kappa-1$, hence
$m_S\ge\kappa$ and the main-box payment applies.
\end{itemize}
This pays every source-free half-range comparison after the last new mode.
Together with the NEW terminal-source edges and the rectangular cases,
the induction is complete, proving (W).

\paragraph{Meaning and formal scope.}
This theorem excludes positive ordinary source recovery throughout the
strict post-LAST-mode reverse first half. It does not assert coefficient
half LC, full property C, the other half of coupled S, or original
one-descent unimodality. The source signs through half follow directly
from (5), including the positive correction one at $M$ in the equal-weight
case, so no unproved coefficient-tilt assertion is hidden in the proof.
The unbounded interval-layer and induction assembly is written here.
Generic scoped kernels are identified separately; no end-to-end Lean
formalization of actual arm membership or this theorem is claimed.

\paragraph{Verification and independent review.}
The four reports in
\texttt{agent\_half\_lc\_weak\_window\_transport.lean} check the
crossing-window curvature certificate, finite chaining of adjacent
source comparisons, removal of an old nonnegative-source guard, and
the earlier sufficient gap arithmetic. The four reports in
\texttt{agent\_half\_lc\_interval\_payments.lean} check the exact
integer interval-overlap formula's whole-layer/missing-first-point
payment, its one-step boundary loss, and the actual late/virtual index
payments. The nine reports in \texttt{HalfTailClosureArithmetic.lean}
check complementary source budgets, the three exceptional heads, their
mass/normalization resources, and the source-free/new-source arithmetic.
All seventeen reports pass with only standard kernel axioms and no
\texttt{sorryAx}; both agent logs have zero warnings or errors.
These are scoped generic kernels. The actual-arm identifications,
height-layer expansion and induction assembly remain the unbounded
written proof above. The completed both-nonincreasing first-half theorem
is reused with its previously declared formal scope.

An independent mathematical review checked every induction boundary,
mode transfer, selected height layer, source support and arithmetic case
without finding a gap. Its independent reverse-append diagnostic
\texttt{agent\_weak\_T\_independent\_diagnostic.py/.json} checks all
735078 unordered pairs from the 1212 actual normalized arms of weight
at most eighteen, including rectangles, plus 1600 wide random pairs:
1495817 strict-post-LAST-mode half-range triples all satisfy (W).
The separate source-prefix/assembly QA
\texttt{agent\_half\_lc\_weak\_window\_transport.py/.json} checks
4000 actual append contexts, 700792 nonnegative old-source prefix
coefficients, 120993 upper-half source-monotonicity edges, 54868
transported prefix-tail edges, and 27455 complete new half-tail edges.
These finite checks are diagnostics only; none supplies an unbounded
step in the proof.
